% !TeX root = ncaa-manuscript-en.tex
\section{Experimental Setup}\label{sec:experimental-setup}

This section defines two evaluation tracks: constraint compliance and preferred
candidate selection by the model, and checking the specifications used to
construct learning inputs. It defines the tasks and reference labels, the
metrics and aggregation units, and the training and comparison protocols.
All learning experiments use synthetic candidate selection; the protocol of
the preliminary Alpamayo-R1 diagnostics is given in Supplementary Section S1.

\subsection{Tasks and Data}\label{sec:setup-data}

Each example contains an image, the original CoC, a constraint sentence when
provided by the condition, descriptions of the supplied candidates, and a selection
instruction (Fig.~\ref{fig:method-inputs}). Candidate measurements are supplied
in text; the VLM does not estimate them from the image. Static tasks have two
candidates; pedestrian entry, maneuver tasks, and position holding have four.

\noindent\textbf{Six behavioral tasks.}
The six tasks form three families. Static tasks apply maximum-speed,
minimum-clearance, and minimum-entry-delay rules to numerical candidate
summaries. Pedestrian-entry examples vary the clearance time, required wait,
candidate entry times, and label assignments. Maneuver tasks concern stopping
and lane changing. Examples occur in pairs with the same scene and candidates
but different supplied rules, so the preferred candidate can change: entry 0.5 s
after a zone clears is admissible under a 0.5 s waiting rule but prohibited
under a 1.5 s rule.

\noindent\textbf{Position holding.}
This task uses the position-hold constraint in Section~\ref{sec:method-eblc}. Synthetic
candidate profiles are presented as XY coordinates or as interval-maximum
distances calculated from the same coordinates; images are diagrams of the
profiles. Each profile yields eight rows: two label rotations, two hold
durations, and two representations. The training, validation, and test splits
contain 120, 16, and 32 profiles, with no profile shared across splits.

Table~\ref{tab:setup-dataset-sizes} gives rows per input condition. P0, P1, and
P2 share examples and splits, differing only in their constraint text. Rows are
constructed examples rather than independent road videos. Pedestrian-entry
images do not overlap across splits, whereas all 13 static primary-test images
also occur in training and maneuver data reuse two diagrams. Static and
maneuver results therefore measure the use of supplied constraints rather than
visual generalization. Split inventories are given in Supplementary Section S2.

\begin{table}[tbp]
\caption{Synthetic dataset sizes per input condition. A dash denotes no
separate validation split.}\label{tab:setup-dataset-sizes}
\centering
\small
\begin{tabular*}{\textwidth}{@{\extracolsep\fill}lrrr@{}}
\toprule
Task family & Training & Validation & Primary test\\
\midrule
Pedestrian entry & 960 & 240 & 240\\
Static (three tasks) & 576 & -- & 144\\
Maneuver (two tasks) & 384 & -- & 96\\
Position hold & 960 & 128 & 256\\
\botrule
\end{tabular*}
\end{table}

\noindent\textbf{Preferred candidates and reference labels.}
As defined in Section~\ref{sec:method-problem}, the preferred candidate has the
highest supplied progress score among admissible candidates; its label is the
reference label. Admissibility is computed from the
supplied rule and candidate values by integer calculations implemented
separately from the logical encoding; position holding checks
$x_i^2+y_i^2\le 200^2$ (mm) for every sample up to the hold endpoint. SMT
candidate verdicts are compared with these calculations
(Section~\ref{sec:results-checks}) and kept as metadata rather than placed in
the prompt. The
supervised response contains only the reference label. Progress scores are
supplied synthetic values whose ordering among admissible candidates
determines the preferred candidate. Labels rotate within each task's label
set: A--B for static tasks and A--D for the other tasks (Supplementary Section S2). Score variation and label rotation reduce fixed
associations between the answer and a label, position, or score, but do not
remove shortcuts based on the supplied numerical values.

\subsection{Metrics and Aggregation}\label{sec:setup-metrics}

\noindent\textbf{Preferred-selection score (PSS)} assigns one to a prediction
that selects the reference label of the highest-scoring admissible candidate,
and zero to every other output. The mean is reported as a percentage. It
measures how often the preferred candidate is selected, not the size of the
progress-score difference when another candidate is selected.
\textbf{Constraint violation rate (CVR)} is the percentage of all predictions
that select a prohibited candidate. An admissible candidate with a lower
progress score receives zero for preferred selection but is not a violation.
The two metrics therefore need not sum to 100\%.
\textbf{Rule-pair accuracy (RPA)} is the percentage of same-scene,
same-candidate pairs with different supplied rules for which both reference
labels are selected. Some position-hold pairs share a reference label, so
this metric does not require a changed choice in every pair.
\textbf{Constraint-compliant completion} is the percentage completing the task without
selecting a prohibited candidate; it is auxiliary. Unextractable answers
receive zero preferred-selection credit and count as neither violations nor
admissible choices; their rate is reported separately. Thus, the complement of the
constraint violation rate is not the admissible-choice rate when extraction
failures occur. Position-hold violations use the strict 0.2 m radius, with a
0.01 m tolerance reported only as an auxiliary measure.

Reported seed standard deviations use the sample convention ($n-1$).
Fresh-adapter results report means and sample standard deviations
over three seeds. Each condition combines both durations, yielding 720
pedestrian-entry predictions (240 rows $\times$ three seeds) or 384
position-hold predictions per representation (128 rows $\times$ three seeds).
Continuation results report counts per duration and
representation (192 predictions from 64 rows and three seeds). Six-task and
additional evaluations report three-seed means. Supplementary Section S6 gives
individual seeds for the fresh-adapter, six-task, representation, and robustness
comparisons; reasoning-prompt and exploratory diagnostics are detailed in
Supplementary Sections S8 and S10. Seeds repeat optimization and do not add
independent examples.

Specification checks comprise the fault-injection evaluation of
Section~\ref{sec:method-checks} (384 cases, 7,392 queries), the 5,376
position-hold candidate checks, and, for the pedestrian-entry specification,
one satisfiability query and 24 boundary-candidate queries for each of four
waits (0.5, 1.5, 0.8, and 1.8 s), 100 queries in total. These units are
reported separately from each other and from model-prediction denominators.
\subsection{Model and Training Settings}\label{sec:setup-training}

We fine-tune Qwen3-VL-2B-Instruct~\cite{qwen2025qwen3vl2b}, a model small
enough for every comparison to be repeated with fresh adapters over three seeds
under matched budgets; a larger model of the same family is evaluated without
training (Section~\ref{sec:setup-comparisons}). Adaptation uses
LoRA~\cite{r14} (rank 8, $\alpha=16$, dropout 0.05) with AdamW (learning rate
$2\times10^{-4}$, weight decay 0.01), batch size 2, and gradient accumulation
over eight steps. Each fit runs for three epochs with optimization seed 42, 17,
or 123; pedestrian-entry and position-hold fits perform 180 updates, and static
and maneuver fits 108 and 72. Within each task family, compared conditions
share examples, epochs, optimizer settings, and the final-checkpoint policy;
validation scores do not select checkpoints. Input lengths differ across text
conditions. Fine-tuned adapters use greedy decoding with a 12-token answer
limit.

\subsection{Comparison Designs}\label{sec:setup-comparisons}

Table~\ref{tab:setup-comparisons} distinguishes three studies by the information
available during training and evaluation. Models trained with constraint text
also receive it at evaluation, matching the intended use of case-specific
constraints (Section~\ref{sec:discussion-training}).

\begin{table}[tbp]
\caption{Training and evaluation information in the three studies. Fresh
adapters share their initialization within each task and seed.}
\label{tab:setup-comparisons}
\footnotesize
\begin{tabularx}{\textwidth}{@{}>{\raggedright\arraybackslash}p{0.14\textwidth}>{\raggedright\arraybackslash}X>{\raggedright\arraybackslash}X>{\raggedright\arraybackslash}X@{}}
\toprule
Study & Training & Evaluation & Question\\
\midrule
Fresh-adapter study & Same pretrained base; CoC alone or CoC + CNL, each with or
without semantic loss & (a) Training input type; (b) same CoC + CNL for all
adapters & With evaluation information and initialization held fixed, do
training constraints reduce violations and improve rule-dependent selection?\\[4pt]
Continuation study & Existing pedestrian-entry adapters; P0: CoC,
P1: CoC + natural-language guard, P2: CoC + CNL & Same CoC + P2 CNL for every
adapter, before and after continuation & Does a constraint-bearing training
history improve behavior under identical evaluation inputs?\\[4pt]
Six tasks & P0: CoC; P1 and P2: CoC + respective constraints & Same input type
as training & Do models use supplied constraints and change their choices when
the rule changes?\\
\botrule
\end{tabularx}
\end{table}

\noindent\textbf{Fresh-adapter study.}
We train new LoRA adapters from the same pretrained base on pedestrian entry and
position holding. The $2\times2$ design crosses CoC-only versus CoC + CNL inputs
with cross-entropy (CE) alone versus CE plus semantic loss; seed-specific
initial LoRA tensors are identical across the four conditions. Two tasks, four
conditions, and three seeds yield 24 fits. Our semantic loss is a categorical
adaptation of Xu et al.~\cite{Xu2018}. At the first supervised answer-token
position, let $z_a$ be the logit for label $a\in\mathcal{Y}=\{A,B,C,D\}$:
\begin{equation}\label{eq:setup-semantic-loss}
q_a=\frac{\exp z_a}{\sum_{b\in\mathcal{Y}}\exp z_b},\qquad
\mathcal{L}_{\mathrm{train}}=\mathcal{L}_{\mathrm{CE}}
-\lambda\log\!\sum_{a\in\mathcal{Y}_C}q_a,
\end{equation}
where $\mathcal{Y}_C\subseteq\mathcal{Y}$ is the set of labels whose candidates
are admissible. We fix $\lambda=1$ before execution. Even with CoC-only inputs,
the loss uses the example's admissible set, so rule information enters through
the objective; inference neither filters candidates nor uses SMT. The protocol
specifies evaluation with the training input type as primary and common-CNL
evaluation as additional; we report both and emphasize the latter when
comparing training conditions.

\noindent\textbf{Continuation study.}
We continue the nine pedestrian-entry adapters (P0, P1, P2 $\times$ three
seeds) on 960 position-hold rows for 180 further updates. P0 receives the
original CoC at both stages, while P1 and P2 receive their respective
constraints at both stages. Before and after continuation, all adapters
receive identical inputs with the P2 constraint sentence. The comparison
therefore includes the effects of different training histories. Previously
inspected profiles with prohibited candidates only 1 mm beyond the radius are
reevaluated separately as a regression check (Supplementary Section S11).

\noindent\textbf{Identifiability without rule text.}
When the distinguishing rule is omitted, identical inputs can have different
reference labels. Every six-task rule pair has different reference labels, giving a
50\% preferred-selection score ceiling for rule-free inputs; in the position-hold set, half the
pairs do, giving an empirical ceiling of 75\%. Fewer violations without rule
text therefore do not show that a model changes its choice with the rule, and
the six-task study measures the use of constraints supplied at both training
and evaluation rather than a training benefit under identical information.
Rule-pair accuracy and the common-CNL evaluation are used for this reason to
assess rule-dependent selection directly.

\noindent\textbf{Pretrained-VLM evaluation.}
Qwen3-VL-2B-Instruct and Qwen3-VL-8B-Instruct~\cite{qwen2025qwen3vl2b,qwen2025qwen3vl8b}
receive CoC and CNL without fine-tuning or demonstrations on all 736
primary-test inputs, with one greedy response per input. The primary evaluation
fixes a 12-token output limit before execution; after observing truncated
explanations, we add a post-hoc evaluation with a 512-token limit and
final-answer extraction, and report both, together with output-limit hits and
answer-extraction failures. This comparison includes adaptation
to the task format and is not a matched-budget training comparison.

\subsection{Comparison Conditions and Related Approaches}\label{sec:setup-related}

The systems reviewed in Section~\ref{sec:related-work} use different
interfaces: DriveLM and DriveVLM address driving question answering and
planning, while nl2spec and NL2TL output temporal-logic formulas rather than
action labels~\cite{r2,r3,r8,r9}. Their native execution paths are documented
in Supplementary Section S5; adding our constraints or an output adapter would
evaluate modified methods. We instead use controls that share selected
principles with prior work (Table~\ref{tab:setup-related-controls}); they are
not reproductions of the cited systems.

\begin{table}[tbp]
\caption{Comparison conditions and principles shared with related approaches.
Correspondences identify selected principles, not reproductions of complete
systems.}\label{tab:setup-related-controls}
\footnotesize
\begin{tabularx}{\textwidth}{@{}>{\raggedright\arraybackslash}X>{\raggedright\arraybackslash}X>{\raggedright\arraybackslash}p{0.25\textwidth}@{}}
\toprule
Condition & Shared principle & Budget\\
\midrule
CNL at evaluation for CoC-only-trained adapters & Use of rules at inference;
regulation-augmented reasoning also uses rules at this stage~\cite{r7} &
Matched fresh-adapter training; continuation reported separately\\[4pt]
CoC + CNL for pretrained 8B & Prompting a larger model with rules & No
task-specific training\\[4pt]
Categorical semantic loss & Constraints in the training objective~\cite{Xu2018}
& Matched fresh-adapter training\\[4pt]
General reasoning prompts & Zero-shot CoT and Plan-and-Solve~\cite{r16,r17}
& No additional training; unequal generation limits\\[4pt]
Logical-text control & Formal representation of requirements~\cite{r8,r9}
& Matched pedestrian-entry training\\[4pt]
P1 and alternative natural-language guards & Natural-language constraints in
training~\cite{r4} & Matched within each task family\\
\botrule
\end{tabularx}
\end{table}

\subsection{Additional Evaluations}\label{sec:setup-additional}

The following evaluations use archived checkpoints, including those from
additional conditions trained with matched task-family budgets. Their rows are
not added to the primary-test denominator. Procedures are detailed in
Supplementary Sections S2, S8, and S10.
\begin{itemize}
\item \textbf{General reasoning prompts.} Zero-shot CoT (128 reasoning tokens)
and basic Plan-and-Solve (256 tokens)~\cite{r16,r17} are applied to the final
P0 checkpoints without rule information or further training. A first stage
generates reasoning for each unique input (546 per method), and a second stage
extracts the label; each method yields 1,440 predictions over three seeds.
\item \textbf{Paraphrasing.} A held-out constraint wording is given to P1 and
P2 on the 240, 144, and 96 primary-test rows of the three families; P0
inputs, which contain no constraint text, are unchanged. A smaller exploratory
diagnostic tests a second wording with 24 rows per wording and seed.
\item \textbf{Unseen values.} Values and durations absent from training are
used, such as pedestrian-entry waits of 0.8 and 1.8 s instead of 0.5 and 1.5 s.
\item \textbf{Combined difficulty.} A 96-row maneuver set combines negation,
clause order, equality boundaries, and m/cm units.
\item \textbf{Alternative pedestrian-entry representations.} We report two
additional conditions trained with the same budget as P0--P2: an alternative
natural-language guard and a logical formula supplied as text. Each was evaluated on the 240 primary-test
and 240 unseen-wait rows. The logical text is a model input, not a solver or
translator output.
\item \textbf{Image interventions.} Pedestrian-entry adapters receive original,
blank, or redrawn images with clearance 1.0 s later, with all other inputs and
the original references fixed (24 rows per condition, 648 predictions). Scores
on redrawn images measure an input intervention rather than scene-understanding performance;
we also report the fraction of choices that change relative to the original
image.
\end{itemize}

\subsection{Bounded Constraint-Composition Diagnostic}\label{sec:setup-composition}

This diagnostic, run after and independently of the training runs, evaluates the manually formalized
source groups in Section~\ref{sec:method-composition}. A source catalog records
the document version, URL, paragraph or page, supported interpretation, and
scope. Eight official documents were inspected; two further UNECE candidates
were excluded because their full official text could not be retrieved.
The executed cases use four document families, with scenario states,
candidate requests, motion assumptions, and test objectives marked as authored
inputs. No operational log is used.

Parameter grids bracket each feasibility threshold by 1 mm: braking (48
cases; speeds of 10/15/20 m/s, decelerations of 5/7 m/s$^2$, and delays of
0/1000 ms, with zero delay as a control), yellow-interval feasibility (64
cases; nominal deceleration of 3 m/s$^2$, not an emergency-braking bound),
all 16 interlock states with drive/hold candidates (32 cases), and all eight
exception states with cross/hold candidates and faithful/omitted-exception
encodings (32 cases). Supplementary Section S13 gives the full grids.

The 176 configurations yield 896 queries, covering the context alone and
every nonempty source subset. Verdicts are compared with independent
rational feasibility calculations, and each SAT witness is checked numerically.
For each UNSAT combination, all proper subsets identify subset-minimal
conflict groups. We distinguish incompatibility introduced by suppression
from avoidance already infeasible without it, and permission to cross from
an obligation to cross. These configurations and queries are counted
separately from fault injection, VLM predictions, and the earlier 123-case
synthetic boundary diagnostic. Supplementary Section S13 records both
protocols and their distinct scopes.

